%! Author = grays
%! Date = 4/3/26

% Preamble
\documentclass[11pt]{article}

% Packages
\usepackage{amsmath}
\usepackage{booktabs}
\usepackage{caption}
\usepackage{float}
\usepackage{graphicx}
\usepackage{listings}
\usepackage{mhchem}
\usepackage{siunitx}
\usepackage{xcolor}
\usepackage[backend=biber,style=ieee]{biblatex}
\usepackage[colorlinks=true, linkcolor=blue, citecolor=blue, urlcolor=blue]{hyperref}
% Listings settings
\lstset{
    basicstyle=\ttfamily\footnotesize,
    keywordstyle=\color{blue},
    commentstyle=\color{green!50!black},
    stringstyle=\color{orange},
    numbers=left,
    numberstyle=\tiny\color{gray},
    stepnumber=1,
    numbersep=5pt,
    backgroundcolor=\color{gray!10},
    showstringspaces=false,
    breaklines=true,
    frame=single,
    rulecolor=\color{gray!50},
    title=\lstname
}
\sisetup{separate-uncertainty=true}
\DeclareSIUnit\angstrom{\text {Å}}

\addbibresource{references.bib}

\title{Franck-Hertz Experiment}
\date{April 3, 2026}
\author{
    Patrick Geraghty (300349204) \\
    Evan Champagne (300443446)
}

% Document
\begin{document}
    \maketitle

    \section{Introduction}\label{sec:intro}

    The Franck-Hertz experiment provides the first direct experimental
    evidence for the quantization of atomic energy levels, confirming
    theoretical predictions initially made by the Bohr model.
    By determining the exact energy required to excite a specific atomic
    transition in mercury, the experimental can be used to calculate an
    experimental value for Planck's constant, \( h \).

    When free electrons are accelerated through mercury vapour, they undergo
    collisions with the mercury atoms.
    If the kinetic energy of an incident electron is less than the energy
    difference between the atom's ground state and its lowest accessible
    excited state (\( \Delta E \)), the collisions are perfectly elastic.
    In this regime, there is no reduction in the electron's kinetic energy.
    However, if the electron's kinetic energy is greater than or equal to
    \( \Delta E \), inelastic collisions can occur.
    During an inelastic collision, the electron transfers a discrete quantum
    of kinetic energy to the mercury atom, exciting it to a higher energy
    state, and continues its path with the surplus energy.
    By observing the energy lost by electrons undergoing these inelastic
    collisions, the excitation energy \( \Delta E \) of the atom can be
    accurately measured.
    This energy difference relates to the accelerating voltage (\( V \)) by
    \begin{equation}
        \Delta E = eV
        \label{eq:energy_diff}
    \end{equation}
    where \( e \) is the elementary charge.

    The specific atomic transition measured in this setup corresponds to the
    excitation of mercury's \( 6s6p^3 P_1 \) state.
    Unlike lower energy transitions, this specific state decays rapidly back
    to the ground state by emitting an ultraviolet photon with a wavelength of
    \( \lambda = 253.7 \) \si{\nano\meter}.
    Using the measured excitation energy and the known wavelength of the
    emitted photon, Planck's constant can be derived using the de Broglie
    relation for photon energy:
    \begin{equation}
        \Delta E = hf = \frac{hc}{\lambda}
        \label{eq:de_broglie_energy}
    \end{equation}
    We can then substitute \( \Delta E \) from Equation~\ref{eq:energy_diff}
    into Equation~\ref{eq:de_broglie_energy} and rearrange to solve for
    \( h \):
    \begin{equation}
        h = \frac{\Delta E \lambda}{c} = \frac{eV \lambda}{c}
        \label{eq:plancks_constant}
    \end{equation}

    \section{Methods}\label{sec:methods}

    \subsection*{Experimental Setup}
    The experimental setup centers around a completely evacuated glass tube
    containing a drop of mercury, which is enclosed within a thermostatically
    controlled oven to precisely regulate the mercury vapour pressure.
    The internal components of the tube operate similarly to a conventional
    vacuum tube, with a filament, a cathode, a grid-like anode, and a
    collector.
    The filament is powered by a 6.3 VAC power supply, which heats the cathode
    to emit electrons via thermionic emission.
    These electrons are then accelerated toward the grounded anode by a
    variable 0--50 V DC power supply.
    Because the anode is grounded, the filament and cathode are configured to
    float with respect to ground, which is necessary for the current
    measurement.
    A picoammeter, capable of reading currents between \SI{10e-9}{\ampere} and
    \SI{10e-7}{\ampere}, is connected to measure the fraction of electrons
    that reach the collector.
    A \SI{1.5}{\volt} battery is placed between the anode and the collector to
    provide a fixed negative retarding potential; this ensures that electrons
    arriving at the anode with less than \SI{1.5}{\electronvolt} of kinetic
    energy are turned back.
    The collector connections are completely shielded to prevent externally
    induced currents, and the battery is carefully insulated from ground to
    stop the \SI{1.5}{\volt} potential from driving an unwanted leakage
    current through the sensitive picoammeter.
    The ambient temperature of the mercury vapour is monitored and adjusted
    using an OMEGA CN370 temperature controller, which reads data from a
    thermocouple attached to a metal ring around the glass tube.

    \subsection*{Experimental Procedure}
    To begin the experiment, the initial temperature on the OMEGA CN370
    controller is set approximately \SI{5}{\degreeCelsius} below the actual
    target temperature before turning on the power switch.
    The optimal operating temperature to observe the maximum number of current
    peaks is typically established between \SI{150}{\degreeCelsius} and
    \SI{180}{\degreeCelsius}, and care is taken to ensure the system never
    exceeds the safety limit of \SI{250}{\degreeCelsius}.
    Once the temperature stabilizes, the anode voltage is gradually increased
    from zero while observing the picoammeter, which must display a current in
    the negative direction.
    Initially, the accelerating voltage is kept below \SI{20}{\volt} to safely
    locate the current peaks without overloading the tube.
    Pushing the voltage too high risks forming a plasma discharge inside the
    tube.
    If a plasma discharge occurs, it acts as a low resistance pathway, causing
    the \SI{1.5}{\volt} battery to drive a large positive current through the
    picoammeter; should this happen, the procedure requires the anode voltage
    to be immediately turned down to zero before resuming.
    By systematically varying the accelerating voltage and recording the
    corresponding collector current, at least four distinct peaks are
    identified.
    This process yields the necessary data to plot the collector current
    against the anode voltage, allowing for the precise measurement of the
    voltage separation between successive peaks.

    \subsection*{Data Analysis}
    The data is collected using an Agilent DSO1012A oscilloscope, which
    captures the collector current signal as a function of time during a
    linear voltage sweep of the anode from 0 to \SI{39}{\volt}.

\printbibliography

\end{document}